\documentclass[a4paper,fleqn]{cas-sc}
% Elsevier CAS single-column template (cas-sc.cls, cas-common.sty from els-cas-templates). Numbered references.
% If the front matter runs over a page, add the longmktitle option.

\usepackage[numbers,sort&compress]{natbib}
\usepackage{amsmath,amssymb}
\usepackage{booktabs}
\usepackage{array}
\usepackage{float}
\usepackage{xurl}

\newtheorem{proposition}{Proposition}
\newtheorem{lemma}{Lemma}
\newdefinition{definition}{Definition}
\newdefinition{remark}{Remark}
\newproof{pf}{Proof}

% Figures are floats. The CAS class reads the optional argument of the figure environment as a key list and defaults
% the placement to top-only, which pushed figures to the end; [pos=htbp] restores the usual placement.

\newcommand{\runhead}[1]{\par\medskip\noindent\textbf{#1}\par\nopagebreak\noindent\ignorespaces}


\newcommand{\figfile}[2]{%
  \IfFileExists{#2.pdf}{\includegraphics[#1]{#2.pdf}}{%
    \IfFileExists{#2.png}{\includegraphics[#1]{#2.png}}{%
      \fbox{\parbox[c][4cm][c]{0.9\linewidth}{\centering\ttfamily\detokenize{#2} (figure file not found)}}}}}

\newcommand{\inputif}[1]{\IfFileExists{#1.tex}{\input{#1.tex}}{%
  \begin{center}\fbox{\ttfamily\detokenize{#1}.tex not found: run bench\_cost.py}\end{center}}}



\begin{document}
\let\WriteBookmarks\relax
\def\floatpagepagefraction{1}
\def\textpagefraction{.001}

\shorttitle{Weak-form rescaling of the initial-condition penalty}
\shortauthors{A. Gouda et al.}

\title [mode = title]{How the weak-form discretization rescales the initial-condition penalty in physics-informed training}

\author[1]{Ayush Gouda}[orcid=0009-0004-4775-3103]
\cormark[1]
\ead{ayushgouda.005@gmail.com}
\credit{Conceptualization, Methodology, Software, Investigation, Formal analysis, Writing -- original draft}

\author[1]{Swaraag {Hebbar N}}
\credit{Writing -- review \& editing}

\author[1]{Aditya Prakash}
\credit{Writing -- review \& editing}

\author[1]{Sudheendra {B S}}
\credit{Supervision, Writing -- review \& editing}

\affiliation[1]{organization={RV Institute of Technology and Management (RVITM)},
            city={Bengaluru},
            country={India}}

\cortext[1]{Corresponding author}

\begin{abstract}
Physics-informed solvers for time-dependent problems often build the Dirichlet data into the trial space and impose the
initial condition through a penalty. This paper gives a scaling law for the weight of that penalty in variational
(weak-in-space) losses and a verification test that exposes its effect, with incompatible initial and boundary data as
the stress case. When the data clash at $t=0$, no smooth trial function meets both, and training on the physics loss
moves a supervised fit to the exact solution away from it: a strong-form loss fails early in the transient and a
weak-form loss late. In the experiments, where the model fails is set by the effective weight of the penalty. The same nominal weight acts
$10^3$ to $10^5$ times more strongly in the weak loss; we call this factor $\kappa$. A scaling lemma shows that on fine meshes $\kappa$
equals the number of test functions per element times $E^{2n}$, with $E$ elements per dimension in $n$ dimensions, and its
error tracks the measured variation of the residual across an element. In 1D, weight sweeps and changes of mesh
and test-function order move each loss's failure as $\kappa$ predicts, on heat and advection--diffusion problems with
exact solutions; in 2D the sweeps meet within a factor of about two of $\kappa$. Normalizing the weak residual by element measure removes the mesh dependence. Loss weights therefore do
not transfer between loss forms or between meshes, and under a clash the effective weight decides which part of the
solution is lost.
\end{abstract}

\begin{highlights}
\item Conflicting boundary and initial data push both losses away from the solution
\item The weak loss weights the initial condition $10^3$ to $10^5$ times more strongly
\item On fine meshes the ratio is test functions per element times $E^{2n}$
\item Weight sweeps and mesh changes move each loss's failure as the ratio predicts
\item A warm-start test exposes the failure while the training loss falls
\end{highlights}

\begin{keywords}
Variational physics-informed neural networks \sep Weak formulation \sep Incompatible initial and boundary data \sep
Loss weighting \sep Heat equation \sep Crank--Nicolson
\end{keywords}


\maketitle

\section{Introduction}
\label{sec:intro}

Physics-informed neural networks place the governing equation in the loss function \cite{raissi2019pinn}. Variational
versions project the residual onto a finite set of test functions and so work with a weak form of the equation
\cite{kharazmi2019vpinn,kharazmi2021hpvpinn,fastvpinn2024}. In both, Dirichlet data can be imposed exactly by building
the trial space around a distance function \cite{lu2021hardconstraint,sukumar2022exact}. Exact imposition was the most
accurate of four Dirichlet strategies compared for PINNs and VPINNs \cite{berrone2023dirichlet}. For a time-dependent
problem, the initial condition is then usually left to the loss as a penalty. This split is close to a default. Such models are used where a classical solver is not a direct substitute, for
example in inverse problems, parametric surrogates and digital twins that combine measured data with physics.

Sometimes the data disagree. Consider a plate at temperature $0$ whose edge is suddenly held at temperature $1$. At the
edge, at $t=0$, the boundary requires $1$ and the initial condition requires $0$. Such incompatible data are a classical
topic for parabolic equations. The solution is not smooth at the corner, and higher regularity needs compatibility
conditions \cite{ladyzhenskaya1968}. Classical discretizations have known treatments. Error estimates for Galerkin
methods with nonsmooth initial data degrade as $t\to0$ \cite{thomee2006}, and Crank--Nicolson time stepping needs a few
damping steps at the start to recover its order \cite{luskin1982,rannacher1984}. A neural trial space with an exact
boundary has a different problem. It is smooth and pinned to the boundary value, so it cannot represent the jump. The
penalty can be driven toward zero only through an ever thinner layer at the edge, which a network of fixed size cannot
follow. Over the models the network can reach, the penalty and the physics residual cannot both be made small, and the
penalty keeps pulling.

Competing loss terms are a known source of failure in physics-informed training
\cite{wang2021gradientflow,wang2022ntk,krishnapriyan2021failure}, and reweighting is the standard response. These studies
balance terms that can all be satisfied at once. Causal training \cite{wang2024causality} reweights the residual in
time so that early times are fitted before later ones; it acts on the same time axis on which the present failure
appears. Corner conflicts are known to hurt training, and the usual remedy is to
drop points near the discontinuity from the loss \cite{stiffreview}. For weak losses, Rojas et al.\ \cite{rojas2024rvpinn}
showed that the plain variational loss need not track the error and proposed a dual-norm loss that does for a rich
enough test space. Robust variational losses for time-dependent problems impose the initial condition through the trial
space, which assumes compatible data \cite{sikora2026crvpinn}. A time-discrete variational scheme was applied to a
freezing problem whose boundary and initial temperatures differ \cite{timediscrete2026}, a case of the same clash in
practice.

This paper gives two tools for this setting: a verification test and a scaling law for loss weights in variational
discretizations. The test is a code-verification test in the sense of \cite{roache1998,oberkampf2010}: a known solution is supplied
and the method is checked against it. It fits a network to the exact solution by supervision and then trains it on the
physics loss alone. Some drift is expected whenever the exact solution lies outside the trial space; the test measures how large it is
and where in the transient it falls. Under the clash both loss forms drift, at opposite ends of the transient, while the
training loss falls. Where the failure falls is set by the effective weight of the initial-condition penalty. The
weak-form discretization multiplies this weight by a factor $\kappa$, which a scaling lemma predicts on fine meshes and
which weight sweeps, changes of mesh and test-function order, and a second equation confirm. The contributions are as
follows.
\begin{itemize}
\item A warm-start verification test that measures how far, and where in time, training leaves a supervised fit
(Section~\ref{sec:default}).
\item A scaling law for the ratio $\kappa$ of the effective weights in the two losses, $\kappa=ME^{2n}$ on fine meshes, with
a measured criterion for when a mesh is fine enough (Section~\ref{sec:kappascaling}).
\item Interventions on the weight, the mesh, the test-function order and the residual normalization, on two equations,
all of which move the failure as $\kappa$ predicts (Sections~\ref{sec:sweeps}--\ref{sec:advdiff}).
\item A classical baseline that places the problem in the trial space, and a comparison of remedies
(Sections~\ref{sec:baseline} and~\ref{sec:remedies}).
\end{itemize}

Section~\ref{sec:formulation} states the problems, the trial space and the losses. Section~\ref{sec:analysis} defines
the effective weight and gives the proposition and the scaling lemma. Section~\ref{sec:design} describes the experiments,
Section~\ref{sec:results} reports the results and Section~\ref{sec:remedies} the remedies. Section~\ref{sec:discussion}
discusses the results, and Section~\ref{sec:conclusion} concludes.

% =============================================================================================
\section{Problem formulation and discretization}
\label{sec:formulation}

\subsection{Model problems}
\label{sec:problems}

The heat equation is solved on $\Omega=(0,1)^n$, $n\in\{1,2\}$, for $t\in(0,1]$:
\begin{equation}
\partial_t u = a\,\Delta u \ \text{in } \Omega, \qquad u = 1 \ \text{on } \partial\Omega, \qquad u(\cdot,0) = u_0,
\label{eq:heat}
\end{equation}
with $a=0.1$. Two initial conditions are used. The \emph{matching} one, $u_0=1-\prod_i\sin(\pi x_i)$, equals $1$ on
$\partial\Omega$, and the exact solution is $u=1-e^{-na\pi^2 t}\prod_i\sin(\pi x_i)$. The \emph{clashing} one is $u_0=0$.
Its exact solution is a Fourier sine series with $400$ odd modes per dimension, accurate for $t\ge10^{-3}$.

The second problem is the advection--diffusion equation
\begin{equation}
\partial_t u + c\,\partial_x u = a\,\partial_{xx} u \ \text{on } (0,1), \qquad u(0,t)=u(1,t)=1,
\label{eq:advdiff}
\end{equation}
with $c=1$, $a=0.1$ and the same two initial conditions. The substitution $u=1-e^{cx/(2a)-c^2t/(4a)}\,w$ turns
Eq.~\eqref{eq:advdiff} into $w_t=a\,w_{xx}$ with $w=0$ at both ends and $w(x,0)=e^{-cx/(2a)}\bigl(1-u_0(x)\bigr)$,
which is solved with a sine series of $799$ modes. This reference has a PDE residual of about $10^{-14}$ and agrees with
an independent Crank--Nicolson solve to about $10^{-6}$.

\subsection{Trial space with an exact boundary}
\label{sec:trial}

The trial functions are
\begin{equation}
u_\theta(x,t) = 1 + B(x)\,N_\theta(x,t), \qquad B(x) = \prod_{i=1}^n 4x_i(1-x_i),
\label{eq:trial}
\end{equation}
where $N_\theta$ is a fully connected network with parameters $\theta$: input $n+1$, five hidden layers of width $64$,
SiLU activations, double precision. Every $u_\theta$ equals $1$ on $\partial\Omega$ for every $t$.

\begin{definition}[Clash]
The initial and boundary data \emph{clash} if $u_0\neq1$ on $\partial\Omega$. Since $u_\theta(\cdot,0)$ is continuous and
equals $1$ on $\partial\Omega$, no trial function then meets the initial condition, and the continuous penalty
$\int_\Omega\bigl(u_\theta(x,0)-u_0(x)\bigr)^2dx$ is positive for every $\theta$.
\end{definition}

The initial-condition penalty is estimated on a set $S$ of points that is redrawn at every step (Eq.~\eqref{eq:ic}), so the
expected value of $L_{\mathrm{IC}}$ is this continuous penalty; a network
cannot drive it to zero by fitting a fixed set of points. Over networks of growing size, the infimum of the loss is nevertheless
zero. The time-shifted solution $u(x,t+\varepsilon)$ is smooth on $\bar\Omega\times[0,1]$, equals $1$ on $\partial\Omega$,
satisfies Eq.~\eqref{eq:heat} exactly, and has an initial-condition error that vanishes as $\varepsilon\to0$. Since
$u(x,t+\varepsilon)-1=B(x)\,w_\varepsilon(x,t)$ with $w_\varepsilon$ smooth, and networks with a smooth non-polynomial
activation such as SiLU approximate smooth functions together with their derivatives on compact sets \cite{pinkus1999},
both terms can be made as small as desired. The infimum is not attained: the shifted solution has a boundary layer of width of order $\sqrt{a\varepsilon}$, thinner the smaller the
error. A network of fixed size trained for a fixed number of steps cannot follow that layer, so over the models it
reaches, the two terms cannot both be made small. Proposition~\ref{prop:linear}(b) describes the resulting trade-off
for a fixed trial space.

\subsection{Losses}
\label{sec:losses}

The loss is
\begin{equation}
L(\theta) = \lambda_{\mathrm{phys}}\,L_{\mathrm{phys}}(\theta) + \lambda_{\mathrm{IC}}\,L_{\mathrm{IC}}(\theta),
\qquad
L_{\mathrm{IC}}(\theta) = \frac{1}{|S|}\sum_{z\in S}\bigl(u_\theta(z,0)-u_0(z)\bigr)^2,
\label{eq:ic}
\end{equation}
with $S$ a set of $500$ points drawn uniformly in $\Omega$ and redrawn at every step, $\lambda_{\mathrm{phys}}=12$ and $\lambda_{\mathrm{IC}}=3$. Two
physics terms are compared.

\runhead{Strong form}
The pointwise residual is averaged over a set $Q$ of space--time points:
\begin{equation}
L_{\mathrm{s}}(\theta) = \frac{1}{|Q|}\sum_{(x,t)\in Q}\bigl(\partial_t u_\theta - a\,\Delta u_\theta\bigr)^2,
\label{eq:strong}
\end{equation}
with $|Q|=1{,}280$ in 1D and $16{,}000$ in 2D, drawn uniformly in $\Omega\times(0,1)$ and redrawn at every step.

\runhead{Weak form}
The domain is split into uniform elements $K_e$, with $E=16$ elements in 1D and $8\times8$ in 2D.
On each element, test functions $v_{e,i}$ are built from $P_{m+1}-P_{m-1}$, with $P_m$ the Legendre polynomials mapped
to $K_e$, and from tensor products of these in 2D. The order $p$ means $m=1,\dots,p+1$, so there are $M=(p+1)^n$ test functions per element and the highest
polynomial degree is $p+2$. The defaults are $p=6$ in 1D ($M=7$) and $p=4$ in 2D ($M=25$). Each test function vanishes
on $\partial K_e$.
Multiplying Eq.~\eqref{eq:heat} by $v_{e,i}$ and applying Green's first identity on $K_e$ gives the residual
\begin{equation}
R_{e,i}(t_j) = \int_{K_e}\bigl(\partial_t u_\theta\, v_{e,i} + a\,\nabla u_\theta\cdot\nabla v_{e,i}\bigr)\,dx,
\label{eq:weakres}
\end{equation}
evaluated by Gauss--Legendre quadrature ($8$ points per element in 1D, $5\times5$ in 2D, raised to $p+2$ per
dimension for higher orders) at $10$ times $t_j$ drawn uniformly in $(0,1)$ at every step. The weak residual needs only
first derivatives of $u_\theta$; the strong residual needs second derivatives. On the default meshes the two losses
evaluate the network at the same number of points per step. The weak loss is the mean square,
\begin{equation}
L_{\mathrm{w}}(\theta) = \frac{1}{N_R}\sum_{e,i,j} R_{e,i}(t_j)^2,
\label{eq:weak}
\end{equation}
with $N_R$ the number of terms. Each time is squared separately; there is no sum over time before squaring. The loss is therefore weak in space
and collocated in time: $\partial_t u_\theta$ enters pointwise at the sampled times, and there are no test functions in time. For
advection--diffusion, $c\,\partial_x u_\theta\,v_{e,i}$ is added to the integrand of Eq.~\eqref{eq:weakres} and
$c\,\partial_x u_\theta$ to the strong residual.

Two variants of the weak loss are also used. The \emph{dual-norm} loss \cite{rojas2024rvpinn} replaces the mean of
squares at each time by $R(t_j)^{\top}G^{-1}R(t_j)$, with $G$ the $H^1$-seminorm Gram matrix of the same test
functions. Its condition number is about $4$ in 1D and $7$ in 2D. The \emph{normalized} loss divides each $R_{e,i}$ by
the element measure $|K_e|$ before squaring. With $E$ uniform elements in 1D, this multiplies $L_{\mathrm{w}}$ by $E^2$.

Neither loss uses gradient clipping. As in FastVPINNs \cite[Eq.~9]{fastvpinn2024}, the cell-wise residual of
Eq.~\eqref{eq:weakres} is an integral over the element and is not divided by the element measure. Codes differ in how
they combine the cell-wise residuals. Eq.~(10) of \cite{fastvpinn2024} averages the squared residual vector over elements,
while its Algorithms~2 and~3, and the hp-VPINN implementation it describes (its Algorithm~1), average over the test
functions of each element and sum over elements. Eq.~\eqref{eq:weak} averages over both.

\subsection{Error measure}
\label{sec:error}

The error is the relative error of the transient,
\begin{equation}
e(t) = \frac{\|u_\theta(\cdot,t)-u(\cdot,t)\|_2}{\|u(\cdot,t)-1\|_2},
\label{eq:error}
\end{equation}
with discrete $L^2$ norms on $201$ points in 1D and $41\times41$ points in 2D, at $t\in\{0.01,0.05,0.1,0.3,1\}$. A model
that outputs the steady state $u=1$ scores $1$. The denominator shrinks as the transient decays, which makes late
errors look large. Forms are therefore compared only at the same time. Absolute errors, $e(t)$ times the RMS of $u-1$,
are given in Table~\ref{tab:abs}.

% =============================================================================================
\section{The effective weight}
\label{sec:analysis}

\subsection{Definition}
\label{sec:effweight}

Let $g_{\mathrm{IC}}=\nabla_\theta\bigl(\lambda_{\mathrm{IC}}L_{\mathrm{IC}}\bigr)$ and
$g_{\mathrm{phys}}=\nabla_\theta\bigl(\lambda_{\mathrm{phys}}L_{\mathrm{phys}}\bigr)$, each averaged over $D=100$ fixed
draws of the sampled points. The $100$-draw estimate is within $4\%$ of a $1{,}000$-draw one (supplementary material). The \emph{effective weight} of the initial-condition term is
\begin{equation}
\eta = \frac{\|g_{\mathrm{IC}}\|}{\|g_{\mathrm{phys}}\|},
\label{eq:eta}
\end{equation}
measured at the supervised fit of Section~\ref{sec:warm}. The initial-condition gradient is the same in both losses,
since it is the same term on the same points. The physics gradient is not. The ratio between the forms is
\begin{equation}
\kappa = \frac{\|g_{\mathrm{phys}}\|_{\mathrm{strong}}}{\|g_{\mathrm{phys}}\|_{\mathrm{weak}}}
= \frac{\eta_{\mathrm{weak}}}{\eta_{\mathrm{strong}}}.
\label{eq:kappa}
\end{equation}
Since $\eta$ is linear in $\lambda_{\mathrm{IC}}$, the weak loss at a given initial-condition weight has the same
effective weight as the strong loss at $\kappa$ times that weight. Adam rescales the gradient
parameter by parameter, so a ratio of global norms does not fix the update on its own. When the weak physics gradient is the strong one
divided by $\kappa$, the total weak-loss gradient is $1/\kappa$ times the strong-loss gradient with the
initial-condition weight multiplied by $\kappa$. Adam is invariant to such an overall rescaling when the gradient is large
compared with its stabilizing constant \cite{kingma2015adam}, so it then takes the same steps for both. The two physics
gradients are close to parallel at the supervised fits (cosine at least $0.987$, Section~\ref{sec:kappascaling}), and
the sweeps of Section~\ref{sec:sweeps} test the resulting prediction directly.

\subsection{A linear trial space}
\label{sec:prop}

For a trial space that is linear in its parameters, the role of $\kappa$ can be stated exactly. This is also the
setting of least-squares training of the last layer of a VPINN \cite{uriarte2025}. Let
$u_\theta = 1 + B(x)\,\Phi(x,t)\,\theta$ with a fixed feature map $\Phi$. Every residual is then affine in $\theta$.
Collect the physics residuals (pointwise values for the strong form, the $R_{e,i}(t_j)$ for the weak form) as
$A\theta-b$ and the initial-condition residuals as $C\theta-d$, with the averaging factors absorbed into $A$ and $C$.
The loss of Eq.~\eqref{eq:ic} becomes
\begin{equation}
L(\theta) = \lambda_{\mathrm{phys}}\,\|A\theta-b\|^2 + \lambda_{\mathrm{IC}}\,\|C\theta-d\|^2 .
\label{eq:linloss}
\end{equation}
Assume $\ker A\cap\ker C=\{0\}$. In this setting the data clash when no $\theta$ satisfies both $A\theta=b$ and
$C\theta=d$.

\begin{proposition}
\label{prop:linear}
Let $\rho=\lambda_{\mathrm{IC}}/\lambda_{\mathrm{phys}}>0$.
\begin{enumerate}
\item[(a)] $L$ has the unique minimizer
\begin{equation}
\theta^*(\rho) = \bigl(A^{\top}A+\rho\,C^{\top}C\bigr)^{-1}\bigl(A^{\top}b+\rho\,C^{\top}d\bigr),
\label{eq:closedform}
\end{equation}
which depends on the weights only through $\rho$.
\item[(b)] The physics term $P(\rho)=\|A\theta^*(\rho)-b\|^2$ is nondecreasing in $\rho$, and the initial-condition term
$I(\rho)=\|C\theta^*(\rho)-d\|^2$ is nonincreasing. Under a clash, $P(\rho)$ and $I(\rho)$ are never both zero.
\item[(c)] If the weak and strong residuals are proportional in the mean-square sense,
$A_{\mathrm{w}}^{\top}A_{\mathrm{w}}=s^2A_{\mathrm{s}}^{\top}A_{\mathrm{s}}$ and
$A_{\mathrm{w}}^{\top}b_{\mathrm{w}}=s^2A_{\mathrm{s}}^{\top}b_{\mathrm{s}}$ for a scalar $s>0$, then $\kappa=1/s^2$ at every $\theta$, and
$\theta^*_{\mathrm{w}}(\rho)=\theta^*_{\mathrm{s}}(\kappa\rho)$.
\end{enumerate}
\end{proposition}

\begin{pf}
(a) Dividing $L$ by $\lambda_{\mathrm{phys}}$ does not move its minimizer, so only $\rho$ matters. The Hessian of
$\|A\theta-b\|^2+\rho\|C\theta-d\|^2$ is $2(A^{\top}A+\rho\,C^{\top}C)$, which is positive definite because
$\ker A\cap\ker C=\{0\}$. Setting the gradient to zero gives Eq.~\eqref{eq:closedform}.

(b) Take $\rho_1<\rho_2$ with minimizers $\theta_1,\theta_2$, and write $P_k=P(\rho_k)$, $I_k=I(\rho_k)$. Optimality of
each minimizer gives $P_1+\rho_1I_1\le P_2+\rho_1I_2$ and $P_2+\rho_2I_2\le P_1+\rho_2I_1$. Adding the two
inequalities gives $(\rho_2-\rho_1)(I_2-I_1)\le0$, so $I_2\le I_1$. The first inequality then gives
$P_1-P_2\le\rho_1(I_2-I_1)\le0$, so $P_1\le P_2$. If $P(\rho)=I(\rho)=0$ for some $\rho$, then $\theta^*(\rho)$
satisfies $A\theta=b$ and $C\theta=d$, which a clash excludes.

(c) The two residual vectors may have different lengths; only their normal equations enter. The physics gradient of
the weak loss is $2\lambda_{\mathrm{phys}}(A_{\mathrm{w}}^{\top}A_{\mathrm{w}}\theta-A_{\mathrm{w}}^{\top}b_{\mathrm{w}})
=s^2\cdot2\lambda_{\mathrm{phys}}(A_{\mathrm{s}}^{\top}A_{\mathrm{s}}\theta-A_{\mathrm{s}}^{\top}b_{\mathrm{s}})$, and the
initial-condition gradient is the same in both losses, so $\kappa=1/s^2$. Substituting into Eq.~\eqref{eq:closedform}
gives $\theta^*_{\mathrm{w}}(\rho)=(s^2A_{\mathrm{s}}^{\top}A_{\mathrm{s}}+\rho\,C^{\top}C)^{-1}
(s^2A_{\mathrm{s}}^{\top}b_{\mathrm{s}}+\rho\,C^{\top}d)=\theta^*_{\mathrm{s}}(\rho/s^2)=\theta^*_{\mathrm{s}}(\kappa\rho)$.
\end{pf}

Parts (a) and (b) are standard results for weighted least squares and scalarized two-objective problems, the
property behind the L-curve of Tikhonov regularization \cite{hansen1992,boyd2004}; part (c) is a rescaling identity.
Part (b) uses only optimality, so it holds for the global minimizers of any loss with two weighted terms: raising the
weight trades physics error for initial-condition error and never improves both. Part (c) is the comparison tested in
Section~\ref{sec:sweeps}: the weak loss behaves like the strong loss with its initial-condition weight multiplied by
$\kappa$.

\begin{remark}
\label{rem:limits}
The proportionality assumed in part (c), and the size of $\kappa$, follow from Lemma~\ref{lem:kappa} for residuals that
are smooth on the element scale. The step from loss terms to the error early and late in the transient is made by the
experiments of Section~\ref{sec:results}, for the network trial space.
\end{remark}

\subsection{Check on a linear model}
\label{sec:lincheck}

The proposition was checked on the clashing 1D heat problem with a linear trial space of $140$ Legendre features in
$(x,t)$ and the exact boundary of Eq.~\eqref{eq:trial}.
The closed form and part (c) reproduce a direct least-squares solve to rounding error, and for the strong, weak and
dual-norm losses $P$ rises and $I$ falls monotonically with $\rho$. For the actual weak residual the linear model has
$\kappa_{\mathrm{lin}}\approx1{,}090$, and the weight that minimizes its error differs between the two losses by a factor
of $100$ to $1{,}000$, the same order. Polynomial features represent the corner layer well, so in this model the two
terms barely conflict and the error stays small. The drift is a property of the network trial space and is studied
with it in Section~\ref{sec:results}.

\subsection{The size of $\kappa$}
\label{sec:kappascaling}

The size of $\kappa$ follows from the structure of the weak loss when the residual is smooth on the element scale.

\begin{lemma}
\label{lem:kappa}
Let the weak loss of Eq.~\eqref{eq:weak} use $E$ uniform elements of width $h=1/E$ per dimension in $n$ dimensions, with
the $M=(p+1)^n$ test functions per element of Section~\ref{sec:losses}, and let the element integrals be exact. Let
$r=\partial_t u_\theta-a\Delta u_\theta$, let $\bar r_e$ be its value at the centre of $K_e$ and $\varrho_e$ its root mean
square over $K_e$, and let $\ell$ be a length such that on every element and at every sampled time
$|\partial_x^{\alpha}r|\le C\,\varrho_e\,\ell^{-|\alpha|}$ for $1\le|\alpha|\le2$, with the same bound for each component of
$\nabla_\theta r$ relative to its own element root mean square. Assume that both losses sample the same distribution of
times, and let $\bar L$ be the space--time mean of $r^2$, which the strong loss of Eq.~\eqref{eq:strong} estimates. Then
\begin{equation}
L_{\mathrm{w}} = \frac{h^{2n}}{M}\,\bar L\,\Bigl(1+O\bigl((h/\ell)^2\bigr)+O\bigl(M(h/\ell)^4\bigr)\Bigr), \qquad
\nabla_\theta L_{\mathrm{w}} = \frac{h^{2n}}{M}\Bigl(\nabla_\theta\bar L+\varepsilon\Bigr),
\label{eq:lemma}
\end{equation}
with $\|\varepsilon\|\le C'\bigl((h/\ell)^2+M(h/\ell)^4\bigr)G$, where $G$ is the mean over elements and times of
$\varrho_e\gamma_e$ and $\gamma_e$ is the element root mean square of $\|\nabla_\theta r\|$. If $\|\nabla_\theta\bar L\|$ is
not small compared with $G$, then, up to the sampling error of the strong loss,
$\kappa = M\,E^{2n}\,\bigl(1+O((h/\ell)^2)+O(M(h/\ell)^4)\bigr)$.
\end{lemma}

\begin{pf}
Since every $v_{e,i}$ vanishes on $\partial K_e$, Green's first identity turns Eq.~\eqref{eq:weakres} into
$R_{e,i}=\int_{K_e} r\,v_{e,i}\,dx$. Take $n=1$ first. Map $K_e$ to $[-1,1]$ by $x=x_e+h\xi/2$, so that
$v_{e,i}(x)=\varphi_m(\xi)$ with $\varphi_m=P_{m+1}-P_{m-1}$, and expand $r=\bar r_e+\tfrac h2 r'_e\,\xi+q_e(\xi)$ about
the element centre, where $|\tfrac h2 r'_e|\le C\varrho_e(h/\ell)$ and $|q_e|\le C\varrho_e(h/\ell)^2$. Since
$\int_{-1}^{1}P_k\,d\xi=2\delta_{k0}$, only $\varphi_1=P_2-P_0$ has a nonzero mean, $\int\varphi_1\,d\xi=-2$, and by parity
$\int\xi\varphi_1\,d\xi=0$, so $R_{e,1}=-h\,\bar r_e+h\,O\bigl(\varrho_e(h/\ell)^2\bigr)$. For $m=2$ the constant term
integrates to zero and $|R_{e,2}|\le Ch\varrho_e(h/\ell)$. For $m\ge3$, $\varphi_m$ is orthogonal to $P_0$ and $P_1$, so only
$q_e$ contributes, and since $|P_k|\le1$ on $[-1,1]$, $|R_{e,m}|\le Ch\varrho_e(h/\ell)^2$ uniformly in $m$. Averaging the
squares over the $M$ test functions, the elements and the sampled times gives
$L_{\mathrm{w}}=(h^2/M)\bigl[\,\overline{\bar r_e^{\,2}}+O\bigl(\overline{\varrho_e^{2}}\,((h/\ell)^2+M(h/\ell)^4)\bigr)\bigr]$.
By the same expansion $\varrho_e^2-\bar r_e^{\,2}=O\bigl(\varrho_e^2(h/\ell)^2\bigr)$ on every element, and the mean of
$\varrho_e^2$ over elements and times is $\bar L$, which gives the first statement. For the gradient,
$\nabla_\theta L_{\mathrm{w}}$ is the mean of $2R_{e,i}\nabla_\theta R_{e,i}$ with
$\nabla_\theta R_{e,i}=\int_{K_e}\nabla_\theta r\,v_{e,i}\,dx$, and the same expansion, applied to $r$ and to
$\nabla_\theta r$, gives the same orders relative to $h^2\varrho_e\gamma_e$, term by term. These bounds are not relative to
the sum, which is why the gradient statement is additive: $\nabla_\theta\bar L$ can be small through cancellation while
each term is not. In two dimensions the test functions are products $\varphi_m(\xi)\varphi_{m'}(\eta)$. Only $(1,1)$ has a
nonzero mean, and the Jacobian $h^2/4$ gives $R_{e,(1,1)}=h^2\,\bar r_e+h^2O(\varrho_e(h/\ell)^2)$; the two products
$(2,1)$ and $(1,2)$ are $O(h^2\varrho_e(h/\ell))$, and all others are $O(h^2\varrho_e(h/\ell)^2)$. The same orders follow
with the factor $h^4/M$. When $\|\nabla_\theta\bar L\|$ is not small compared with $G$, the ratio of gradient norms is
$Mh^{-2n}$ times $1+O((h/\ell)^2)+O(M(h/\ell)^4)$.
\end{pf}

On a fine mesh the weak residual is, to leading order, the element-centre residual scaled by $h^n$ and carried by one
test function per element. The term $M(h/\ell)^4$ comes from the test functions of order three and higher, each of which
sees only the curvature of the residual; it is no larger than the first correction while $(h/\ell)^2\le1/M$. This is the proportionality assumed in Proposition~\ref{prop:linear}(c), with $s^2=h^{2n}/M$.
The factor $h^{2n}$ comes from integrating over small elements without dividing by their measure, and $1/M$ from
averaging over $M$ test functions, of which a smooth residual excites one. Both are fixed by the discretization. The same calculation gives $\kappa$ for the other
conventions of Section~\ref{sec:losses}: $\kappa=E^{2n}$ when the squared residuals are averaged over elements and summed
over test functions (Eq.~10 of \cite{fastvpinn2024}), and $\kappa=ME^{n}$ when they are averaged over test functions and
summed over elements (Algorithms~2 and~3 of \cite{fastvpinn2024}, and the hp-VPINN implementation they describe).
Two codes with the same mesh and test space therefore differ in $\kappa$ by a factor of $M$ or $E^n$, and a weight tuned in one
does not carry to the other.

Table~\ref{tab:kappa} and Figure~\ref{fig:kappa} compare the prediction with $\kappa$ measured at $100$ draws. With the
mesh fixed at $16$ elements, $\kappa$ is $1.21$ to $1.25$ times $256(p+1)$ for every order, so the order dependence is
the predicted one; the linear model of Section~\ref{sec:lincheck} shows the same slope. On $16$ and $32$ elements the
prediction is within $25\%$ for both initial conditions, both equations and warm and cold starts. On coarser meshes
$\kappa$ is larger than $ME^{2n}$, by $4.4$ on $4$ elements, $2.9$ on $8$ and $2.1$ on the $8\times8$ mesh in 2D, and the
slope fitted over $4$ to $32$ elements is $1.24$, rising to $1.69$ between $16$ and $32$. The lemma is an asymptotic
statement and needs $h$ small compared with $\ell$, the length over which the residual changes. On the default meshes
$(h/\ell)^2$ is $6.5$ in 1D and $17$ in 2D (Table~\ref{tab:align}), so agreement there is measured rather than implied by
the lemma; the relation below quantifies it.

Table~\ref{tab:align} tests the lemma's hypotheses at the supervised fits. The size of $h/\ell$ is estimated on every
element and sampled time from a Legendre fit of $r$ at the element's quadrature points,
$r\approx c_0+c_1P_1(\xi)+c_2P_2(\xi)$, as $h/\ell_e=\max\bigl(2|c_1|/|c_0|,\sqrt{12|c_2|/|c_0|}\bigr)$ (the tensor
version in 2D), and averaged with weights $c_0^2$. The weak and strong physics gradients point the same way, with cosine
$0.987$ or more on every mesh and $0.9999$ on $32$ elements. The gradient of a mean-square residual is a sum of terms
$r\,\nabla_\theta r$ over points, and the lemma's gradient statement needs these terms not to cancel. We measure this with
$\Gamma=\bigl\|\sum r\,\nabla_\theta r\bigr\|\big/\sum|r|\,\|\nabla_\theta r\|$, taken over random space--time points:
$\Gamma=1$ means the terms all point the same way, and $\Gamma$ near $0$ means they largely cancel. The gradient shows no
strong cancellation ($\Gamma$ of $0.84$ to $0.99$ under the clash). The weighted $(h/\ell)^2$ falls with refinement, from about $106$ on $4$ elements to
$1.3$ on $32$, and the deviation of $\kappa$ from $ME^{2n}$ tracks it: where the deviation is above the noise of a few
percent, $\kappa/(ME^{2n})-1$ is $0.03$ to $0.07$ times $(h/\ell)^2$, on every 1D mesh and in 2D. The constant depends on
the estimator; the relation does not. It gives a working criterion: $(h/\ell)^2$ below about $4$ at a supervised fit or a
reference solution keeps the prediction within roughly $25\%$, and estimating it needs only residual evaluations. The
same relation accounts for the effect of the clash on $\kappa$. The clashing fits have the larger $(h/\ell)^2$ ($6.5$
against $0.4$ in 1D, $17$ against $1.7$ in 2D) and the larger $\kappa$, by $20\%$ in 1D and a factor of $2$ in 2D. Where the residual is smooth, at the random
initializations of the cold starts, $\kappa$ is within $2\%$ of $ME^{2n}$ in both 1D and 2D (Table~\ref{tab:kappa}).

The cosine between the initial-condition and physics gradients measures the clash directly. With matching data it is
positive ($+0.07$ to $+0.71$): both terms pull the same way. Under the clash it is negative ($-0.47$ to $-0.96$), except for
one 2D seed ($+0.63$): the initial-condition term pulls against the physics.

\begin{table}[htbp]
\centering
\caption{Checks of Lemma~\ref{lem:kappa} at the supervised fits, warm start, seeds 42--44 (ranges over seeds; $(h/\ell)^2$:
mean over seeds of the $c_0^2$-weighted mean). Cosines are between the averaged gradients over $D=100$ draws; $\Gamma$ is
the cancellation ratio of the text. Ratio: measured $\kappa$ ($100$ draws, GPU) over $ME^{2n}$; the last column is omitted where the deviation is within the noise. Default mesh: $16$
elements, $p=6$ in 1D; $8\times8$, $p=4$ in 2D.}
\label{tab:align}
\small
\setlength{\tabcolsep}{3pt}
\begin{tabular}{lcccccc}
\toprule
Setting & cos(weak, strong) & cos(IC, weak) & $\Gamma$ & $(h/\ell)^2$ & ratio & $(\text{ratio}-1)/(h/\ell)^2$ \\
\midrule
1D clash, $E=4$ & 0.987--0.994 & $-0.75$ to $-0.40$ & 0.96--0.99 & 106 & 4.4 & 0.032 \\
1D clash, $E=8$ & 0.997--0.999 & $-0.77$ to $-0.46$ & 0.96--0.99 & 40 & 2.9 & 0.047 \\
1D clash, $E=16$ & 0.9991--0.9996 & $-0.79$ to $-0.47$ & 0.96--0.99 & 6.5 & 1.22 & 0.034 \\
1D clash, $E=32$ & 0.9998--0.9999 & $-0.80$ to $-0.47$ & 0.96--0.99 & 1.3 & 0.99 & -- \\
1D matching, $E=16$ & 1.000 & $+0.07$ to $+0.40$ & 0.60--0.65 & 0.38 & 1.02 & -- \\
2D clash, $8\times8$ & 0.999--1.000 & $-0.96$ to $+0.63$ & 0.84--0.86 & 17 & 2.1 & 0.067 \\
2D matching, $8\times8$ & 0.999--1.000 & $+0.46$ to $+0.71$ & 0.41--0.55 & 1.7 & 1.06 & -- \\
\bottomrule
\end{tabular}
\end{table}

\begin{table}[htbp]
\centering
\caption{Measured $\kappa$ ($100$ draws, geometric mean over seeds 42--44, warm start unless stated) against
Lemma~\ref{lem:kappa}, $\kappa\approx M E^{2n}$ with $M=(p+1)^n$. Default: 1D heat, clashing initial condition, $16$
elements, $p=6$ ($M=7$, prediction $1{,}792$). $\kappa_{\mathrm{norm}}=\kappa/E^2$ is the ratio after normalizing the
weak residual by element measure, for which the lemma predicts $M=7$.}
\label{tab:kappa}
\small
\setlength{\tabcolsep}{4pt}
\begin{tabular}{lcccc}
\toprule
Elements $E$ ($p=6$) & 4 & 8 & 16 & 32 \\
$\kappa$ & 490 & 1{,}300 & 2{,}200 & 7{,}100 \\
$7E^2$ & 112 & 448 & 1{,}792 & 7{,}168 \\
ratio & 4.4 & 2.9 & 1.23 & 0.99 \\
$\kappa_{\mathrm{norm}}$ (lemma: $7$) & 30 & 20 & 8.6 & 6.9 \\
\midrule
Order $p$ ($E=16$) & 2 & 4 & 6 & 10 \\
$\kappa$ & 940 & 1{,}600 & 2{,}200 & 3{,}400 \\
$256(p+1)$ & 768 & 1{,}280 & 1{,}792 & 2{,}816 \\
ratio & 1.22 & 1.25 & 1.23 & 1.21 \\
\midrule
Other settings & \multicolumn{2}{c}{$\kappa$} & prediction & ratio \\
1D heat, matching IC & \multicolumn{2}{c}{1{,}810--1{,}820} & 1{,}792 & 1.01--1.02 \\
1D heat, cold start & \multicolumn{2}{c}{1{,}750--1{,}790} & 1{,}792 & 0.98--1.00 \\
1D advection--diffusion & \multicolumn{2}{c}{2{,}000} & 1{,}792 & 1.12 \\
2D heat ($E=8$, $p=4$, $M=25$) & \multicolumn{2}{c}{$2.1$--$2.3\times10^5$} & $1.02\times10^5$ & 2.0--2.2 \\
2D heat, matching IC & \multicolumn{2}{c}{$1.06$--$1.11\times10^5$} & $1.02\times10^5$ & 1.04--1.08 \\
2D heat, cold start & \multicolumn{2}{c}{$1.01$--$1.04\times10^5$} & $1.02\times10^5$ & 0.99--1.02 \\
\bottomrule
\end{tabular}
\end{table}

\begin{figure}[pos=htbp]
\centering
\figfile{width=\linewidth,height=0.42\textheight,keepaspectratio}{fig_kappa}
\caption{The weak-form discretization sets $\kappa$. Left: $\kappa$ against the number of elements $E$ ($p=6$). Right:
$\kappa$ against the test-function order plus one ($E=16$). 1D heat, clashing initial condition, warm start, $100$ draws.
Filled: network (geometric mean over seeds 42--44, bars: min--max). Open: linear model of Section~\ref{sec:lincheck}.
Solid and dashed lines: least-squares fits. Dotted: Lemma~\ref{lem:kappa}, $7E^2$ (left) and $256(p+1)$ (right).}
\label{fig:kappa}
\end{figure}

% =============================================================================================
\section{Experimental design}
\label{sec:design}

\subsection{Warm and cold starts}
\label{sec:warm}

A \emph{warm} start is a network fitted to the exact solution by supervision for $5{,}000$ Adam steps, then trained on
the loss of Eq.~\eqref{eq:ic} for $10{,}000$ steps at learning rate $10^{-4}$. If the loss is minimized at the solution,
training should not make the fit worse. With a matching initial condition it can improve it. The supervised fits are
within about $0.07$ of the exact solution at every time (Table~\ref{tab:giveup}). They are not exact, since the exact
solution cannot be represented under the clash. A \emph{cold} start is a random initialization trained for $15{,}000$
steps at learning rate $10^{-3}$. All runs use cosine decay.

Arms use seeds 42--44. The three 2D weak-form arms at default weights with the clash (plain and dual-norm from warm
starts, plain from cold starts) use seeds 42--51, because their late-time spread is large
(supplementary material).

\subsection{Weight sweeps and the match criterion}
\label{sec:sweepdesign}

The initial-condition weight is multiplied by a fixed factor in one form and compared with the other form at its default
weight. The strong form is swept upward and the weak form downward, from the same warm starts, with seeds 42--44 and
everything else unchanged. The multipliers were fixed in advance and are listed in the supplementary material. Sweep comparisons use the
three-seed means of the default arms as reference, so like is compared with like.

The \emph{gap} between two arms is the largest difference of their mean errors over the five times. Two arms
\emph{match} when the gap is at most $0.05$. This threshold is above the seed range of every 1D weak-form arm at $t=1$
(at most $0.036$). A threshold decides how many points pass, but says little about where they lie, and any monotone
trade-off crosses the other form's default somewhere. The main test is therefore where the gap is smallest, and whether
that is where the independently measured $\kappa$ says.

\subsection{Classical baseline}
\label{sec:cndesign}

The classical baseline solves the same heat problems with second-order central differences in space on uniform grids
whose nodes include every evaluation point, $h=1/(200m)$ in 1D and $h=1/(40m)$ in 2D, with $m\in\{1,2,4,8\}$ in 1D and
$\{1,2,4\}$ in 2D. Time steps are $\Delta t=0.01/q$ with $q\in\{2,4,8,16\}$, so every evaluation time is hit exactly.
Three time schemes are used: plain Crank--Nicolson \cite{cranknicolson1947}, Crank--Nicolson with Rannacher start-up,
in which the first two steps are replaced by four backward Euler half-steps \cite{rannacher1984}, and backward Euler. The
reference and the error measure are those of Section~\ref{sec:problems}.

% =============================================================================================
\section{Results}
\label{sec:results}

\subsection{Both losses leave the supervised fit}
\label{sec:default}

Figure~\ref{fig:heat} shows the warm-start runs at default weights. With a matching initial condition no arm drifts.
With a clashing one both losses drift, at opposite ends of the transient. In 2D the strong form's error is $0.42$ at
$t=0.01$ and falls to $0.007$ at $t=1$, while the weak form's is $0.084$ early and grows to $0.32$ at $t=1$ (median $0.27$
over ten seeds). In 1D the strong form is again worse early ($0.28$ against $0.083$), and the weak form's error is about
twice the supervised fit's throughout. Without the initial-condition term neither form learns the transient (error
$0.36$ to $1.0$).

\begin{figure}[pos=htbp]
\centering
\figfile{width=\linewidth,height=0.42\textheight,keepaspectratio}{fig_heat}
\caption{Heat equation, warm start, default weights: relative error of the transient over time. Lines: mean over seeds;
bands: min--max. The 2D clashing weak-form arms use ten seeds and all other arms three. Dashed grey: the finest
finite-difference solution with Rannacher start-up (Section~\ref{sec:baseline}). With a matching initial condition (left)
no arm drifts. With a clashing one (right) the strong form's error is large early and falls, and the weak form's is
small early and grows. The dual-norm weak loss behaves like the plain one. With the initial condition off, the weak form
stays far from the transient.}
\label{fig:heat}
\end{figure}

Table~\ref{tab:giveup} shows how each loss resolves the conflict in 2D. The strong form gives up the initial condition
(penalty $0.25\to0.92$) and repairs the physics (pointwise residual $38\to0.046$). The weak form keeps the
initial-condition fit and leaves the residual at about $52$.

\begin{table}[htbp]
\centering
\caption{Heat 2D, clashing initial condition, warm start: loss terms and error. Means over $n$ seeds; at $t=1$ also
the median and [min--max]. The pointwise PDE residual is $12$ times the mean of $(\partial_tu-a\Delta u)^2$ over $16{,}000$
fixed random points. The dual-norm residual was scored on a different set of points.}
\label{tab:giveup}
\footnotesize
\setlength{\tabcolsep}{3pt}
\begin{tabular}{lccccc}
\toprule
Model & $n$ & IC loss & PDE residual & $e(0.01)$ & $e(1)$: mean, median [min--max] \\
\midrule
Supervised start & 10 & 0.253 & 37.9 & 0.065 & 0.040 \\
Weak & 10 & 0.233 & 52.2 & 0.084 & 0.320, 0.269 [0.081--0.721] \\
Weak, dual norm & 10 & 0.233 & 53.2 & 0.084 & 0.389, 0.330 [0.118--0.795] \\
Strong & 3 & 0.921 & 0.046 & 0.415 & 0.007, 0.007 [0.006--0.008] \\
\bottomrule
\end{tabular}
\end{table}

\runhead{Absolute errors}
The relative error of Eq.~\eqref{eq:error} divides by the size of the transient, which
shrinks with time. In absolute terms (Table~\ref{tab:abs}) the pattern holds. In 2D the strong form's error peaks at
$t=0.01$ ($0.36$) and the weak form's at $t=0.1$ ($0.13$), and from $t=0.3$ on the weak form's error is $4$ to $44$ times
the strong form's.

\begin{table}[htbp]
\centering
\caption{Heat, clashing initial condition, warm start, default weights: absolute error, the relative error of
Eq.~\eqref{eq:error} times the RMS of $u-1$ on the evaluation grid, mean over $n$ seeds. ``Transient'' is the RMS of
$u-1$ itself.}
\label{tab:abs}
\small
\setlength{\tabcolsep}{4pt}
\begin{tabular}{llccccc}
\toprule
 & & $t=0.01$ & $t=0.05$ & $t=0.1$ & $t=0.3$ & $t=1$ \\
\midrule
1D & Transient & 0.95 & 0.88 & 0.82 & 0.67 & 0.34 \\
 & Strong ($n=3$) & 0.27 & 0.16 & 0.098 & 0.016 & 0.0013 \\
 & Weak ($n=3$) & 0.079 & 0.072 & 0.070 & 0.016 & 0.010 \\
\midrule
2D & Transient & 0.88 & 0.76 & 0.66 & 0.44 & 0.11 \\
 & Strong ($n=3$) & 0.36 & 0.21 & 0.12 & 0.016 & 0.00079 \\
 & Weak ($n=10$) & 0.074 & 0.11 & 0.13 & 0.060 & 0.035 \\
\bottomrule
\end{tabular}
\end{table}

\runhead{The loss prefers the drifted model}
Over weak-form training in 2D both terms of the objective fall, the
initial-condition term from $0.253$ to $0.233$ and the weak residual from $6.6\times10^{-5}$ to $3.9\times10^{-5}$, while
the error at $t=1$ rises eightfold, from $0.040$ to $0.320$. Among the models the network can reach, the loss ranks the
drifted one above the supervised fit. A falling training loss is therefore no evidence that the model approaches the
solution, and the warm-start test detects what the loss does not.

\runhead{Seeds and cold starts}
The late-time 2D weak-form errors vary between seeds ($0.081$ to $0.72$ at $t=1$ over ten
seeds), while the early errors agree closely (standard deviation $0.006$ at $t=0.01$). From random initializations the
same split appears: in 2D the strong form is worse at $t=0.01$ ($0.23$ against $0.13$) and recovers to $0.005$, while the
weak form reaches $0.59$ at $t=1$.

\subsection{The effective weight moves the failure}
\label{sec:sweeps}

$\kappa$ is $2{,}200$ on 1D heat and $2.2\times10^5$ on 2D heat (Table~\ref{tab:kappa}).

\runhead{A large effective weight is harmless without a clash}
With a matching initial condition the weak form's
initial-condition term has an effective weight of $125$ to $637$ in 2D, and nothing drifts. The scale gap is as large as
under the clash; only the clash is removed. The clash causes the failure, and the effective weight decides where it
shows. This separates the effect from loss imbalance as usually studied \cite{wang2021gradientflow,wang2022ntk}.

\runhead{Weight sweeps}
Figure~\ref{fig:sweep} shows the sweeps; the full tables are in the supplementary material.
Raising the strong form's weight moves its error from early to late: in 1D its error at $t=0.01$ falls from $0.28$ to
$0.078$ while its error at $t=1$ rises from $0.004$ to $0.052$. Its gap to the weak default is smallest at the measured
$\kappa$, $0.004$ at $2{,}200$ times its default, and under $0.02$ from $940$ to $3{,}400$ times. Lowering the weak form's
weight does the reverse. Its gap to the strong default is smallest at about $2.2/\kappa$ in both 1D ($0.016$) and 2D
($0.014$), where the two forms also reach the same initial-condition loss ($0.944$ against $0.921$ in 2D). In 1D the two
sweeps trace one curve: a weak run and a strong run whose weights differ by $1{,}000$ to $3{,}000$ agree in their errors
and in their loss terms, for example in initial-condition losses of $0.131$ and $0.125$ and pointwise residuals of $20.1$
and $20.0$. In 2D the strong form takes on the weak form's behaviour from $10^4$ times its default upward, its early
error falling from $0.42$ to $0.094$ and its initial-condition loss reaching the weak form's $0.233$. Its closest point,
$0.55\kappa$, has a worst gap of $0.096$, set by the late-time seed spread of the weak default. A Gaussian-process fit to
each sweep's gap places the minimum between about $0.7$ and $2.3$ times $\kappa$ in all six sweeps (supplementary
material).

\begin{figure}[pos=htbp]
\centering
\figfile{width=\linewidth,height=0.42\textheight,keepaspectratio}{fig_sweep}
\caption{The initial-condition weight moves where each loss fails. Clashing initial condition, warm start, seeds
42--44; left: 1D heat, right: 2D heat. Lines: mean; bands: min--max. Early ($t=0.01$) and late ($t=1$) error against the
initial-condition weight relative to the strong form's default. Weak-form runs are placed at their multiplier times
$\kappa$ ($2{,}200$ in 1D, $2.2\times10^5$ in 2D). Raising the strong form's weight (filled) moves its error from early
to late; lowering the weak form's (open) does the reverse. In 1D the two fall on one curve. Dashed: $\kappa$.}
\label{fig:sweep}
\end{figure}

\subsection{Mesh, order and normalization}
\label{sec:mesh}

\runhead{Mesh}
Changing only the weak form's mesh changes $\kappa$ (Table~\ref{tab:kappa}), and the weak form's failure
moves with it: its error at $t=0.05$ rises from $0.031$ on $4$ elements to $0.120$ on $32$ (Table~\ref{tab:mesh}). The
strong form at $\kappa$ times its default follows it within $0.024$ on every mesh. The same loss at the same nominal
weight fails differently on a different mesh, as $\kappa$ predicts.

\begin{table}[htbp]
\centering
\caption{Weak-form mesh, 1D heat, clashing initial condition, warm start: relative error, mean over seeds 42--44. The
strong form is unchanged except for its initial-condition weight, set to $\kappa$ times its default, with $\kappa$
measured on each mesh. ``Gap'' is the largest difference between the two rows over the five times.}
\label{tab:mesh}
\small
\setlength{\tabcolsep}{4pt}
\begin{tabular}{rrlcccccc}
\toprule
$E$ & $\kappa$ & Arm & $t=0.01$ & $t=0.05$ & $t=0.1$ & $t=0.3$ & $t=1$ & gap \\
\midrule
4 & 490 & weak, default & 0.057 & 0.031 & 0.043 & 0.021 & 0.020 & 0.024 \\
 & & strong, $\times\kappa$ & 0.078 & 0.055 & 0.056 & 0.022 & 0.021 & \\
\addlinespace
8 & 1{,}300 & weak, default & 0.065 & 0.056 & 0.062 & 0.023 & 0.021 & 0.015 \\
 & & strong, $\times\kappa$ & 0.080 & 0.070 & 0.073 & 0.023 & 0.031 & \\
\addlinespace
16 & 2{,}200 & weak, default & 0.083 & 0.082 & 0.085 & 0.024 & 0.031 & 0.004 \\
 & & strong, $\times\kappa$ & 0.083 & 0.080 & 0.083 & 0.021 & 0.035 & \\
\addlinespace
32 & 7{,}100 & weak, default & 0.087 & 0.120 & 0.115 & 0.032 & 0.038 & 0.009 \\
 & & strong, $\times\kappa$ & 0.089 & 0.120 & 0.113 & 0.033 & 0.047 & \\
\bottomrule
\end{tabular}
\end{table}

\runhead{Order}
On $16$ elements $\kappa$ varies by a factor of $3.6$ across orders $p=2$ to $10$, as
Lemma~\ref{lem:kappa} predicts ($11/3\approx3.7$), and at each order the weak form at its default weight lands within
$0.005$ of the strong form at $\kappa(p)$ times its default (Table~\ref{tab:order}).

\runhead{Normalization}
Dividing the weak residual by its element measure multiplies the physics term by $E^2$, and
$\kappa$ falls by that factor, to $\kappa_{\mathrm{norm}}=30$, $20$, $8.6$ and $6.9$ on $4$ to $32$ elements. These values
approach the $M=7$ that Lemma~\ref{lem:kappa} leaves after normalization. Under Adam this is the same as lowering the weak form's initial-condition weight by
$E^2$, and the failure follows the rescaling. The normalized weak form behaves like the strong form at $7$ to $30$ times
its default, its failure moves from late to early (error at $t=1$ of $0.005$ against $0.031$ on $16$ elements), and its
dependence on the mesh shrinks: the largest spread of its error over the four meshes falls from $0.089$ to $0.032$. Over
all eleven settings of mesh, order and normalization, the weak form matches the strong form at the corresponding $\kappa$
within $0.03$.

\begin{table}[htbp]
\centering
\caption{Test-function order and normalization, 1D heat, clashing initial condition, warm start, seeds 42--44. ``Gap''
is the largest difference over the five times from the strong form at $\kappa$ (order) or $\kappa_{\mathrm{norm}}$
(normalized) times its default.}
\label{tab:order}
\small
\begin{tabular}{lcccc}
\toprule
Order $p$ ($16$ elements) & 2 & 4 & 6 & 10 \\
$\kappa$ & 940 & 1{,}600 & 2{,}200 & 3{,}400 \\
gap & 0.005 & 0.003 & 0.004 & 0.005 \\
\midrule
Elements $E$ (normalized) & 4 & 8 & 16 & 32 \\
$\kappa_{\mathrm{norm}}$ & 30 & 20 & 8.6 & 6.9 \\
gap & 0.021 & 0.024 & 0.026 & 0.029 \\
\bottomrule
\end{tabular}
\end{table}

\subsection{A second equation}
\label{sec:advdiff}

The main tests were repeated on the advection--diffusion problem of Eq.~\eqref{eq:advdiff} with the same trial space,
losses and network. With the clash $\kappa=2{,}000$, against $2{,}200$ for 1D heat on the same mesh, so the ratio belongs
to the discretization of the loss rather than to the equation. The defaults split as on heat (Table~\ref{tab:adv}), and
the forms meet near the measured ratio. The strong form's gap to the weak default is smallest at $1.5\kappa$ ($0.014$)
and the weak form's at $2/\kappa$ ($0.007$), where the loss terms also agree (initial-condition loss $0.064$ against
$0.066$, pointwise residual $75.2$ against $75.0$). With a matching initial condition the strong form at $\kappa$ times its
default stays within $0.014$ of its start, so a large weight is again harmless when it can be satisfied.

\begin{table}[htbp]
\centering
\caption{Advection--diffusion, clashing initial condition, warm start, mean over seeds 42--44. Top: relative error at
default weights. Bottom: gap to the other form's default for each initial-condition multiplier. $\kappa=2{,}000$.}
\label{tab:adv}
\small
\setlength{\tabcolsep}{3pt}
\begin{tabular}{lcccccc}
\toprule
Default weights & $t=0.01$ & $t=0.05$ & $t=0.1$ & $t=0.3$ & $t=1$ & \\
\midrule
Strong & 0.171 & 0.087 & 0.064 & 0.050 & 0.056 & \\
Weak & 0.058 & 0.033 & 0.048 & 0.088 & 0.178 & \\
\midrule
Strong multiplier & $\times100$ & $\times300$ & $\times1000$ & $\times2000$ & $\times3000$ & $\times10^4$ \\
Gap to weak default & 0.153 & 0.120 & 0.066 & 0.030 & 0.014 & 0.070 \\
\midrule
Weak multiplier & $\times10^{-4}$ & $\times3{\cdot}10^{-4}$ & $\times5{\cdot}10^{-4}$ & $\times10^{-3}$ & $\times3{\cdot}10^{-3}$ & $\times10^{-2}$ \\
Gap to strong default & 0.280 & 0.119 & 0.053 & 0.007 & 0.044 & 0.090 \\
\bottomrule
\end{tabular}
\end{table}

\subsection{A classical baseline}
\label{sec:baseline}

Table~\ref{tab:cn} shows the same clash for a classical discretization. Plain Crank--Nicolson shows the textbook failure
under incompatible data: its amplification factor for stiff modes is close to $-1$, so the error from the jump at $t=0$
decays slowly \cite{luskin1982}. In 1D at $\Delta t=5\times10^{-3}$, refining the mesh raises its error at $t=0.3$ from
$6.6\times10^{-4}$ to $8.1\times10^{-3}$. Rannacher start-up \cite{rannacher1984} removes this and restores second order in
time, and with matching data plain Crank--Nicolson converges at second order. The finest Rannacher solutions are within
$3.3\times10^{-4}$ of the reference at every time.

The comparison isolates the cause. The grid stores the jump at $t=0$ exactly, with the boundary nodes at $1$ and the
interior nodes at $0$, and Rannacher start-up repairs a time-stepping problem. A smooth trial function pinned to the
boundary can approach the jump only through a layer that a fixed network cannot follow, which is why the two loss
terms pull against each other. The clash is a property of the
trial space and the loss, not of the equation.

\begin{table}[htbp]
\centering
\caption{Clashing initial condition: relative error of a finite-difference solver and of the trained models (warm start,
default weights). FD rows: second-order central differences; ``R'' is Crank--Nicolson with Rannacher start-up, ``CN''
plain Crank--Nicolson. Model rows: mean over $n$ seeds.}
\label{tab:cn}
\footnotesize
\setlength{\tabcolsep}{2.5pt}
\begin{tabular}{llccccc}
\toprule
 & Solver & $t=0.01$ & $t=0.05$ & $t=0.1$ & $t=0.3$ & $t=1$ \\
\midrule
1D & R, $h=1/1600$, $\Delta t=6.25{\cdot}10^{-4}$ & $1.2{\cdot}10^{-4}$ & $7.5{\cdot}10^{-6}$ & $2.5{\cdot}10^{-6}$ & $4.2{\cdot}10^{-7}$ & $1.5{\cdot}10^{-7}$ \\
 & CN, $h=1/1600$, $\Delta t=6.25{\cdot}10^{-4}$ & $2.5{\cdot}10^{-3}$ & $1.9{\cdot}10^{-4}$ & $6.6{\cdot}10^{-5}$ & $1.1{\cdot}10^{-6}$ & $3.6{\cdot}10^{-8}$ \\
 & R, $h=1/1600$, $\Delta t=5{\cdot}10^{-3}$ & $9.7{\cdot}10^{-3}$ & $4.8{\cdot}10^{-4}$ & $1.5{\cdot}10^{-4}$ & $2.1{\cdot}10^{-5}$ & $1.0{\cdot}10^{-5}$ \\
 & CN, $h=1/1600$, $\Delta t=5{\cdot}10^{-3}$ & $5.2{\cdot}10^{-2}$ & $1.6{\cdot}10^{-2}$ & $5.9{\cdot}10^{-3}$ & $8.1{\cdot}10^{-3}$ & $9.7{\cdot}10^{-4}$ \\
 & Strong ($n=3$) & 0.282 & 0.180 & 0.119 & 0.024 & 0.004 \\
 & Weak ($n=3$) & 0.083 & 0.082 & 0.085 & 0.024 & 0.031 \\
\midrule
2D & R, $h=1/160$, $\Delta t=6.25{\cdot}10^{-4}$ & $3.3{\cdot}10^{-4}$ & $9.7{\cdot}10^{-5}$ & $6.6{\cdot}10^{-5}$ & $4.9{\cdot}10^{-5}$ & $5.2{\cdot}10^{-7}$ \\
 & Strong ($n=3$) & 0.415 & 0.273 & 0.182 & 0.037 & 0.007 \\
 & Weak ($n=10$) & 0.084 & 0.150 & 0.192 & 0.136 & 0.320 \\
\bottomrule
\end{tabular}
\end{table}

\subsection{Computational cost}
\label{sec:cost}

Table~\ref{tab:cost} gives the cost per training step of each loss and the cost of measuring $\kappa$, timed on one GPU
with no other jobs running. On the default meshes both losses evaluate the network at the same number of points per step
($1{,}280$ in 1D, $16{,}000$ in 2D). Per step, the strong loss needs a forward pass and $n+1$ backward passes through
the network at its $|Q|$ points, one for the first derivatives and one per dimension for the second. The weak loss needs a
forward pass and one backward pass at its $10E^nQ^n$ quadrature points, with $Q$ points per dimension, followed by a
contraction with the test functions that costs $O(10E^nMQ^n)$ operations and is small next to the network passes. The weak loss is therefore cheaper per step, by a factor of $1.6$ in 1D and $2.1$
in 2D, and in 2D it uses about a third of the memory. A warm-started run takes $2$ to $5$ minutes. Measuring $\kappa$ takes
$D=100$ gradient evaluations of each loss term at a single point in parameter space, $2.7$\,s per seed in 1D and $3.9$\,s in
2D, about $2$ to $3\%$ of a warm weak-form run. On fine meshes Lemma~\ref{lem:kappa} gives $\kappa$ with no measurement at all.

The finite-difference solver of Section~\ref{sec:baseline} takes $0.16$\,s in 1D and about $8$\,s in 2D on a CPU.
Physics-informed models are used where such a solver is not a substitute, and the warm-start test of
Section~\ref{sec:default} can use a classical solve of this kind as its reference.

\inputif{cost_table}

% =============================================================================================
\section{Remedies}
\label{sec:remedies}

\subsection{Intermediate weights}
Since the weight moves the failure between the two ends of the transient, intermediate
weights reduce the worst error. The strong form at $100$ to $300$ times its default keeps its 1D error at or below $0.10$
at every time, against $0.28$ at the default, and the weak form at $10^{-3}$ times its default in 2D stays at or below
$0.11$, against $0.351$ at $t=1$ at the default (three-seed means, as in the sweep). The trade-off itself remains, as
Proposition~\ref{prop:linear}(b) predicts for the loss terms.

\subsection{Learning-rate annealing}
Learning-rate annealing \cite{wang2021gradientflow} resets the initial-condition weight
during training from the ratio of the largest physics gradient to the mean initial-condition gradient
(Table~\ref{tab:anneal}). It sets large weights in every case, $2\times10^4$ to $1.3\times10^5$ times the default for the
1D strong form without the clash and $2\times10^7$ to $5\times10^7$ times with it. Without the clash this does no harm,
and every warm-started arm stays within $0.017$ of its default. With the clash it moves the failure rather than removing
it: the annealed 2D strong form lands next to its sweep point at $1.2\times10^5$ times its default, and in 1D both forms
are pushed beyond the sweeps, to errors near $1$ mid-transient.

\begin{table}[htbp]
\centering
\caption{Learning-rate annealing of the initial-condition weight \cite{wang2021gradientflow}, warm start: relative error,
mean over seeds 42--44. ``Change'' is the largest difference from the same form at its default weight over the five
times.}
\label{tab:anneal}
\small
\setlength{\tabcolsep}{4pt}
\begin{tabular}{lllcccccc}
\toprule
 & IC & Form & $t=0.01$ & $t=0.05$ & $t=0.1$ & $t=0.3$ & $t=1$ & change \\
\midrule
1D & matching & weak & 0.001 & 0.002 & 0.003 & 0.006 & 0.004 & 0.005 \\
 & & strong & 0.002 & 0.002 & 0.003 & 0.005 & 0.003 & 0.005 \\
 & clashing & weak & 0.084 & 0.247 & 0.458 & 0.807 & 0.743 & 0.783 \\
 & & strong & 0.097 & 0.380 & 0.671 & 0.977 & 0.447 & 0.953 \\
\midrule
2D & matching & weak & 0.003 & 0.004 & 0.004 & 0.007 & 0.037 & 0.013 \\
 & & strong & 0.003 & 0.003 & 0.004 & 0.005 & 0.018 & 0.017 \\
 & clashing & weak & 0.087 & 0.107 & 0.102 & 0.048 & 0.179 & 0.172 \\
 & & strong & 0.095 & 0.166 & 0.198 & 0.081 & 0.451 & 0.444 \\
\bottomrule
\end{tabular}
\end{table}

\subsection{Ramping the boundary}
Switching the boundary on smoothly, $g(t)=1-e^{-t/\tau}$, removes the jump at $t=0$. On 1D
heat with $\tau=0.05$ and $0.01$, the weak form and the strong form at $630$ and $870$ times its default (a third to a half
of $\kappa=1{,}900$) stay within $0.05$ of their supervised starts at every time, while the strong form at its default
weight still gives up the initial condition (changes of $0.095$ and $0.146$). A quick ramp leaves a steep layer at the
corner, so the weight still decides the outcome.

\subsection{Dual-norm weighting}
With the element-local test spaces used here, the dual-norm loss \cite{rojas2024rvpinn}
does not remove the drift.
In 2D it is slightly worse than the plain weak loss ($0.389$ against $0.320$ at $t=1$, worse in nine of ten paired seeds,
paired $t$-test $p=0.0007$). Its guarantee holds for test spaces rich enough to represent the residual \cite{rojas2024rvpinn}.

\subsection{Masking the initial condition}
Dropping initial-condition points within a band of the boundary
\cite{stiffreview} reduces the 2D weak form's late error ($0.351\to0.107$ warm, $0.485\to0.262$ cold) but raises the strong
form's ($0.007\to0.066$ warm). Removing the initial condition stops the drift but loses the transient.

\subsection{Normalization}
Normalizing the weak residual by element measure is the remedy for carrying weights between
loss forms and meshes (Section~\ref{sec:mesh}).

% =============================================================================================
\section{Discussion}
\label{sec:discussion}

\subsection{What the effective weight does}
When the two terms cannot both be made small, the initial-condition term keeps
pulling, and the effective weight decides which gives way. When it is small, the model gives up the initial condition
and repairs the physics, so it is wrong early and recovers. When it is large, the model holds the initial condition and
leaves the physics unrepaired, so it drifts late. Proposition~\ref{prop:linear} shows this trade-off for the loss terms
of a linear trial space, and the sweeps show it for the network: the switch is gradual, and annealing, which raises the
weight automatically, moves along it.

\subsection{Weights do not transfer}
The same nominal weight means different things in a weak and a strong loss, and in a
weak loss on different meshes. On fine meshes the factor is $ME^{2n}$ (Lemma~\ref{lem:kappa}) and can be computed before
training; on coarser meshes $(h/\ell)^2$ at a fit shows how far it moves. Normalizing the weak residual by element
measure removes the $E^{2n}$ part. On a non-uniform mesh the same calculation weights each element by $h_e^{2n}$, so the
unnormalized loss favours large elements, and normalization removes that as well. Under a clash this matters more than usual, because the effective weight decides which
part of the solution is lost.

\subsection{A test that can be run}
Fit the network to a reference, train it on the physics loss alone, and see whether its
error grows. On these problems the test exposes the failure in both forms while the training loss falls; on 2D heat the
weak objective falls while the late error grows eightfold. A classical solve, such as the finite-difference solver of
Section~\ref{sec:baseline}, can serve as the reference.

\subsection{Higher dimensions}
The factor grows quickly with dimension. In 3D with $8$ elements per dimension and $p=4$,
the formula of Lemma~\ref{lem:kappa} gives $ME^{2n}=125\cdot8^6\approx3.3\times10^7$. Normalization by element measure removes the $E^{2n}$ part at
no cost.

\subsection{Where the clash arises}
Any model that builds a Dirichlet boundary into the trial space and keeps the initial
condition soft meets the clash when the boundary is switched on suddenly: a wall brought to a new temperature, a step
load, a surface exposed to a new environment at $t=0$. The cheap actions are to smooth the switch-on where the problem
allows, to normalize the weak residual by element measure, to compute or measure $\kappa$ before reusing weights, and to
run the warm-start test when a reference exists.

\subsection{Scope}
The model problems are linear, on the unit square, with exact solutions, so every claim is checked
against the exact answer. One network architecture and one family of test functions are used, the loss is weak in space
and collocated in time, and the mesh, order and normalization interventions are 1D. Training uses Adam. Most arms use three seeds and
the 2D weak-form defaults ten.

% =============================================================================================
\section{Conclusion}
\label{sec:conclusion}

When a hard-coded boundary contradicts a soft initial condition, a network of fixed size cannot make the
initial-condition penalty small without hurting the physics, and its pull moves both weak-form and strong-form training away from a supervised fit to the exact
solution, at opposite ends of the transient. Where the model goes wrong is set by the effective weight of the penalty.
The weak-form discretization multiplies that weight by $\kappa=ME^{2n}$ on fine meshes, with a correction that tracks the
variation of the residual across an element. Weight sweeps and changes of mesh and test-function order move each form's
failure as $\kappa$ predicts, on heat and on advection--diffusion. A finite-difference solver with Rannacher start-up
resolves the same clash, which places the problem in the trial space and the loss. In practice: smooth the switch-on
where the problem allows, normalize the weak residual by element measure, compute or measure $\kappa$ before carrying
weights between loss forms or meshes, and test a physics loss by starting from a supervised fit and checking whether its
error grows.

\printcredits

\section*{Declaration of competing interest}
The authors declare that they have no known competing financial interests or personal relationships that could have
appeared to influence the work reported in this paper.

\section*{Funding}

This research did not receive any specific grant from funding agencies in the public, commercial, or not-for-profit
sectors.

\section*{Data availability}
The code and all logged results, including per-seed values for every number in the text, are available at
\url{https://github.com/AyushG-1210/weakform-ic-penalty}.

\section*{Declaration of generative AI and AI-assisted technologies in the writing process}
During the preparation of this work the authors used AI assistants (Claude, Anthropic; Gemini, Google) to draft and edit
the text, to write and debug code, and to help with analysis. After using these tools, the authors reviewed and edited
the content as needed and take full responsibility for the content of the publication.

\begin{thebibliography}{30}

\bibitem{raissi2019pinn}
M. Raissi, P. Perdikaris, G.E. Karniadakis, Physics-informed neural networks: A deep learning framework for solving
forward and inverse problems involving nonlinear partial differential equations, J. Comput. Phys. 378 (2019) 686--707.

\bibitem{kharazmi2019vpinn}
E. Kharazmi, Z. Zhang, G.E. Karniadakis, Variational physics-informed neural networks for solving partial
differential equations, arXiv preprint arXiv:1912.00873, 2019.

\bibitem{kharazmi2021hpvpinn}
E. Kharazmi, Z. Zhang, G.E. Karniadakis, hp-VPINNs: Variational physics-informed neural networks with domain
decomposition, Comput. Methods Appl. Mech. Engrg. 374 (2021) 113547.

\bibitem{fastvpinn2024}
T. Anandh, D. Ghose, H. Jain, S. Ganesan, FastVPINNs: Tensor-driven acceleration of VPINNs for complex geometries,
arXiv preprint arXiv:2404.12063, 2024.

\bibitem{lu2021hardconstraint}
L. Lu, R. Pestourie, W. Yao, Z. Wang, F. Verdugo, S.G. Johnson, Physics-informed neural networks with hard constraints
for inverse design, SIAM J. Sci. Comput. 43 (6) (2021) B1105--B1132.

\bibitem{sukumar2022exact}
N. Sukumar, A. Srivastava, Exact imposition of boundary conditions with distance functions in physics-informed deep
neural networks, Comput. Methods Appl. Mech. Engrg. 389 (2022) 114333.

\bibitem{berrone2023dirichlet}
S. Berrone, C. Canuto, M. Pintore, Enforcing Dirichlet boundary conditions in physics-informed neural networks and
variational physics-informed neural networks, Heliyon 9 (8) (2023) e18820.

\bibitem{ladyzhenskaya1968}
O.A. Ladyzhenskaya, V.A. Solonnikov, N.N. Ural'tseva, Linear and Quasi-linear Equations of Parabolic Type, Transl.
Math. Monogr., vol. 23, American Mathematical Society, Providence, RI, 1968.

\bibitem{thomee2006}
V. Thom\'ee, Galerkin Finite Element Methods for Parabolic Problems, 2nd ed., Springer Ser. Comput. Math., vol. 25,
Springer, Berlin, 2006.

\bibitem{luskin1982}
M. Luskin, R. Rannacher, On the smoothing property of the Crank--Nicolson scheme, Appl. Anal. 14 (2) (1982) 117--135.

\bibitem{rannacher1984}
R. Rannacher, Finite element solution of diffusion problems with irregular data, Numer. Math. 43 (2) (1984) 309--327.

\bibitem{wang2021gradientflow}
S. Wang, Y. Teng, P. Perdikaris, Understanding and mitigating gradient flow pathologies in physics-informed neural
networks, SIAM J. Sci. Comput. 43 (5) (2021) A3055--A3081.

\bibitem{wang2022ntk}
S. Wang, X. Yu, P. Perdikaris, When and why PINNs fail to train: A neural tangent kernel perspective, J. Comput. Phys.
449 (2022) 110768.

\bibitem{krishnapriyan2021failure}
A. Krishnapriyan, A. Gholami, S. Zhe, R. Kirby, M.W. Mahoney, Characterizing possible failure modes in
physics-informed neural networks, in: Advances in Neural Information Processing Systems 34 (2021) 26548--26560.

\bibitem{wang2024causality}
S. Wang, S. Sankaran, P. Perdikaris, Respecting causality for training physics-informed neural networks, Comput. Methods
Appl. Mech. Engrg. 421 (2024) 116813.

\bibitem{stiffreview}
P. Sharma, L. Evans, M. Tindall, P. Nithiarasu, Stiff-PDEs and physics-informed neural networks, Arch. Comput. Methods
Eng. 30 (5) (2023) 2929--2958.

\bibitem{rojas2024rvpinn}
S. Rojas, P. Maczuga, J. Mu\~noz-Matute, D. Pardo, M. Paszy\'nski, Robust variational physics-informed neural
networks, Comput. Methods Appl. Mech. Engrg. 425 (2024) 116904.

\bibitem{sikora2026crvpinn}
M. Sikora, L. Siwik, N. Leszczy\'nska, T.M. Ciesielski, E. Valseth, M. Bastidas Olivares, M. {\L}o\'s, T.
S{\l}u\.zalec, J. Leszczy\'nski, M. Paszy\'nski, Collocation-based robust physics informed neural networks for
time-dependent simulations of pollution propagation under thermal inversion conditions on Spitsbergen, arXiv preprint
arXiv:2604.23003, 2026.

\bibitem{timediscrete2026}
M. Bastidas Olivares, J.D. Acosta Castrill\'on, D.A. Mu\~noz, Physics-informed deep learning for industrial processes:
time-discrete VPINNs for heat conduction, arXiv preprint arXiv:2603.04711, 2026.

\bibitem{roache1998}
P.J. Roache, Verification and Validation in Computational Science and Engineering, Hermosa Publishers, Albuquerque, NM,
1998.

\bibitem{oberkampf2010}
W.L. Oberkampf, C.J. Roy, Verification and Validation in Scientific Computing, Cambridge University Press, Cambridge,
2010.

\bibitem{pinkus1999}
A. Pinkus, Approximation theory of the MLP model in neural networks, Acta Numer. 8 (1999) 143--195.

\bibitem{kingma2015adam}
D.P. Kingma, J. Ba, Adam: A method for stochastic optimization, in: International Conference on Learning
Representations (ICLR), 2015.

\bibitem{uriarte2025}
C. Uriarte, M. Bastidas, D. Pardo, J.M. Taylor, S. Rojas, Optimizing variational physics-informed neural networks using
least squares, Comput. Math. Appl. 185 (2025) 76--93.

\bibitem{hansen1992}
P.C. Hansen, Analysis of discrete ill-posed problems by means of the L-curve, SIAM Rev. 34 (4) (1992) 561--580.

\bibitem{boyd2004}
S. Boyd, L. Vandenberghe, Convex Optimization, Cambridge University Press, Cambridge, 2004.

\bibitem{cranknicolson1947}
J. Crank, P. Nicolson, A practical method for numerical evaluation of solutions of partial differential equations of
the heat-conduction type, Math. Proc. Cambridge Philos. Soc. 43 (1) (1947) 50--67.

\end{thebibliography}


\clearpage\addtocounter{page}{-1}
\end{document}